% © 2026 Ing. David Jaroš — CC BY-NC-ND 4.0
%
% File: research_tracks/BRIDGE_CLOSURE/B3_weinberg_odd_spinor_bridge_status.tex
% Status: BRIDGE CLOSURE NOTE.

\documentclass[12pt,a4paper]{article}
\usepackage[a4paper,margin=2.5cm]{geometry}
\usepackage{amsmath,amssymb,booktabs}
\usepackage[most]{tcolorbox}
\usepackage{hyperref}
\usepackage{xcolor}
\usepackage{microtype}

\hypersetup{colorlinks=true,linkcolor=blue!60!black,urlcolor=blue!60!black,citecolor=green!50!black}

\newtcolorbox{statusbox}[1][]{
  colback=green!5!white,
  colframe=green!55!black,
  arc=3pt,
  boxrule=1pt,
  title=Weinberg Bridge Status,
  #1
}

\title{Bridge B3: Weinberg Odd-Spinor Threshold Status}
\author{Ing. David Jaroš}
\date{2026-06-14}

\begin{document}
\maketitle

\begin{abstract}
This note records the current Weinberg-angle bridge status.  The GUT boundary
$\sin^2\theta_W=3/8$ is derived from hypercharge norms.  The odd-spinor
threshold spectrum gives $b_1=33/5$, $b_2=1$, $b_3=-3$.  Combining these with
the compact scale bridge and the alpha/QED running bridge gives the $Z$-pole
closure $\sin^2\theta_W(M_Z)=0.23122\ldots$.
\end{abstract}

\begin{statusbox}
Current status:
\[
  \sin^2\theta_W(M_Z)=0.23122\ldots
\]
\[
  \text{PROVEN relative to current UBT bridge assumptions.}
\]
\end{statusbox}

\section{Key files}

\begin{itemize}
  \item \texttt{weinberg\_angle\_gut\_norm\_theorem.tex}
  \item \texttt{weinberg\_angle\_odd\_spinor\_threshold\_beta\_theorem.tex}
  \item \texttt{weinberg\_angle\_mz\_bridge\_closure.tex}
\end{itemize}

\section{Remaining external-facing caveat}

The bridge proof depends on the current UBT interpretation of the odd-spinor
threshold as the active RG spectrum and on the compact-scale calibration.  This
is acceptable inside the repository, but a standalone paper must defend these
bridges explicitly.

\end{document}
